\documentclass[11pt]{article}
\usepackage[a4paper,margin=1in]{geometry}
\usepackage{amsmath,amssymb,amsthm}
\usepackage{booktabs}
\usepackage{enumitem}
\usepackage{hyperref}
\usepackage{microtype}

\newtheorem{theorem}{Theorem}[section]
\newtheorem{lemma}[theorem]{Lemma}
\newtheorem{corollary}[theorem]{Corollary}
\newtheorem{proposition}[theorem]{Proposition}
\theoremstyle{definition}
\newtheorem{definition}[theorem]{Definition}
\theoremstyle{remark}
\newtheorem{remark}[theorem]{Remark}

\newcommand{\NP}{\mathsf{NP}}
\newcommand{\Z}{\mathbb Z}

\title{Closing the Rank-Four Gap for $f$-Factors in Balanced Hypergraphs}
\author{}
\date{}

\begin{document}
\maketitle

\begin{abstract}
We study exact incidence-degree problems in balanced hypergraphs. Previous
work established polynomial solvability for rank at most three and
NP-completeness for non-uniform balanced hypergraphs using a rank-five
construction, leaving rank four open. We close this gap. We give a simple
rank-four selector gadget and prove that it is balanced. A reduction from
bounded-occurrence three-dimensional matching shows that both the
$f$-factor problem and the perfect $f$-matching problem are NP-complete even
when the maximum degree is at most ten and $\max_v f(v)\le5$. Applying the
standard vertex-expansion operation yields NP-completeness of ordinary
Perfect Matching for non-uniform Mengerian hypergraphs of rank at most four
and maximum degree at most fifteen. Since Mengerian hypergraphs are ideal,
the same instances give an ideal-hypergraph corollary. Thus the complexity
dichotomy by rank bound for non-uniform balanced $f$-factor and perfect
$f$-matching problems is polynomial for bounds at most three and
NP-complete already for bound four.
\end{abstract}

\section{Introduction}

Let $H=(V,E)$ be a finite hypergraph. For a function
$f:V\to\Z_{\ge0}$, an $f$-factor is a subset of edges in which vertex $v$
has incidence degree exactly $f(v)$. The related perfect $f$-matching
problem permits nonnegative integer multiplicities on edges. These are
natural exact-covering problems, and their complexity depends strongly on
the incidence structure of $H$.

Balanced hypergraphs are defined by the absence of strong odd cycles. Their
incidence matrices have strong integrality properties; in particular,
balanced hypergraphs of rank at most three have totally unimodular incidence
matrices. Beckenbach and Scheidweiler established polynomial solvability in
that range and proved NP-completeness for non-uniform balanced hypergraphs
with a rank-five reduction \cite{BS17}. Their paper explicitly identified
rank four as the remaining case; the same boundary is recorded in the more
detailed dissertation account \cite{Beckenbach2019}.

We prove that rank four is already hard. The central gadget replaces the
rank-five selector edge used previously. For every source triple, it has a
center $c$, helpers $h_1,\ldots,h_4$, auxiliary vertices $r_1,r_2$, four
fallback edges, two rank-four selector edges, one rank-three synchronizing
edge, and three source ports. Two equations at $r_1,r_2$ synchronize the
selector edges; the center equation then forces either zero or all three
ports. The only delicate point is global balancedness, which follows from a
case analysis of strong cycles.

Our results are as follows.

\begin{theorem}[Balanced rank-four factors]
The $f$-factor problem and the perfect $f$-matching problem are
NP-complete for non-uniform balanced hypergraphs of rank at most four. The
hardness holds for instances satisfying
\[
\Delta(H)\le10,\qquad \max_{v\in V}f(v)\le5.
\]
\end{theorem}

\begin{theorem}[Mengerian rank-four perfect matching]
Perfect Matching is NP-complete for non-uniform Mengerian hypergraphs of
rank at most four and maximum degree at most fifteen.
\end{theorem}

The construction is non-uniform. We make no claim about uniform rank-four
hypergraphs, nor about counting versions of the problems.

\section{Preliminaries}

We use simple hypergraphs: distinct edges are distinct subsets of $V$, and
no repeated edges occur. The degree of $v$ is the number of edges containing
$v$, and $\Delta(H)$ is the maximum degree. The rank is
$\max_{e\in E}|e|$.

\begin{definition}[Strong cycle and balancedness]
A strong cycle of length $k\ge2$ is a cyclic sequence of distinct vertices
$v_1,\ldots,v_k$ and distinct edges $e_1,\ldots,e_k$ such that
\[
e_i\cap\{v_1,\ldots,v_k\}=\{v_i,v_{i+1}\}
\]
for every $i$, with indices modulo $k$. A hypergraph is balanced if it has
no strong cycle of odd length.
\end{definition}

\begin{definition}[$f$-factors and perfect $f$-matchings]
An $f$-factor is a subset $F\subseteq E$ satisfying
\[
|\{e\in F:v\in e\}|=f(v)\quad(v\in V).
\]
A perfect $f$-matching is a vector $x\in\Z_{\ge0}^{E}$ satisfying
\[
\sum_{e\ni v}x(e)=f(v)\quad(v\in V).
\]
When $f\equiv1$, the latter is an ordinary perfect matching.
\end{definition}

\begin{definition}[Mengerian and ideal hypergraphs]
For a hypergraph $K$, let $\nu(K)$ be its maximum matching size and
$\tau(K)$ its minimum vertex-cover size. For a nonnegative integer vector
$a$ on $V(K)$, let $K^a$ denote the vertex expansion obtained by replacing
$v$ with $a(v)$ clones and replacing every incident edge by all choices of
one clone. A hypergraph is Mengerian if
\[
\nu(K^a)=\tau(K^a)
\]
for every $a$. It is ideal if its vertex-cover polyhedron
\[
\{y\in\mathbb R_{\ge0}^{V}:\sum_{v\in e}y_v\ge1\ (e\in E)\}
\]
has only integral vertices.
\end{definition}

Balanced hypergraphs are Mengerian, and every Mengerian hypergraph is ideal;
these are standard min--max consequences for balanced incidence matrices
and their expansions \cite{BS17,Beckenbach2019}.

\section{The rank-four construction}

We reduce from three-dimensional matching. The source consists of disjoint
sets $X,Y,Z$ of equal size and a set
$M\subseteq X\times Y\times Z$. A perfect 3-dimensional matching is a set
of triples covering every source vertex exactly once. We use the standard
bounded-occurrence version in which every source vertex occurs in at most
three triples; this version is NP-complete \cite{GJ79}.

For each $e=(v_1,v_2,v_3)\in M$, introduce private vertices
\[
c_e,h_{e,1},h_{e,2},h_{e,3},h_{e,4},r_{e,1},r_{e,2}.
\]
The edges of the gadget are
\begin{align*}
U_{e,i}&=\{c_e,h_{e,i}\} &&(1\le i\le4),\\
B_{e,1}&=\{c_e,h_{e,1},h_{e,2},r_{e,1}\},\\
B_{e,2}&=\{c_e,h_{e,3},h_{e,4},r_{e,2}\},\\
A_e&=\{c_e,r_{e,1},r_{e,2}\},\\
P_{e,j}&=\{c_e,v_j\} &&(1\le j\le3).
\end{align*}
Set $f(c_e)=5$ and $f(w)=1$ for every other vertex. The only shared
vertices are the source vertices; all gadget vertices are private.

\begin{lemma}[Two local modes]
In every $f$-factor, each gadget is in exactly one of the following modes.
\begin{enumerate}[label=(\roman*)]
\item OFF: $A_e$ and all four $U_{e,i}$ are selected, while no $B$-edge or
source port is selected.
\item ON: both $B$-edges and all three source ports are selected, while
$A_e$ and all four $U$-edges are unselected.
\end{enumerate}
Both modes satisfy all local equations.
\end{lemma}

\begin{proof}
Write $u_i,b_1,b_2,a,p_j\in\{0,1\}$ for the edge indicators. The
$r$-vertices give
\[
b_1+a=1,\qquad b_2+a=1.
\]
Thus $b_1=b_2=b$ and $a=1-b$. Each helper gives $u_i+b=1$, hence
$u_i=1-b$. At the center,
\[
4(1-b)+2b+(1-b)+\sum_{j=1}^3p_j=5,
\]
so $\sum_jp_j=3b$. If $b=0$, all ports are absent; if $b=1$, all three
ports are present. Substitution verifies both modes.
\end{proof}

\section{Balancedness}

The following lemma is the structural core of the reduction.

\begin{lemma}
Every strong cycle of length greater than two in the constructed hypergraph
uses only size-two edges.
\end{lemma}

\begin{proof}
The only edges of size greater than two are $B_{e,1},B_{e,2},A_e$. Fix a
gadget $e$ and suppose first that $c_e$ is a cycle vertex. A non-size-two
edge containing $c_e$ can contain at most one other cycle vertex. If that
vertex is an $h_{e,i}$, the only other edge incident with it is $U_{e,i}$;
the two cycle edges at it therefore return immediately to $c_e$, producing
only a 2-cycle. If the other vertex is $r_{e,1}$ (the other case is
symmetric), continuation through $A_e$ either uses $r_{e,2}$ as well, in
which case $A_e$ contains the three cycle vertices $c_e,r_{e,1},r_{e,2}$,
contrary to strength, or uses $c_e$ and closes a 2-cycle. The same argument
with $B_{e,1}$ applies when the current edge is $A_e$.

Now suppose $c_e$ is not a cycle vertex. No size-two edge incident with
$c_e$ can be a cycle edge, because it contains only one cycle vertex. Thus
no $U$- or $P$-edge occurs. Among the remaining edges, the private supports
intersect only as
\[
B_{e,1}\cap A_e=\{r_{e,1}\},\quad
A_e\cap B_{e,2}=\{r_{e,2}\},\quad
B_{e,1}\cap B_{e,2}=\varnothing.
\]
They form a path, not a cycle. A two-edge strong cycle is impossible because
each pair shares at most one vertex, while a 2-cycle requires two common
cycle vertices. Different gadgets have disjoint private vertices. Hence no
strong cycle avoiding $c_e$ can use a non-size-two edge.

It follows that every strong cycle of length greater than two consists only
of $U$- and $P$-edges. These form a bipartite graph with gadget centers on
one side and source/helper vertices on the other, so every such cycle is
even. Therefore the hypergraph is balanced.
\end{proof}

\section{Hardness for balanced hypergraphs}

\begin{theorem}
The construction has an $f$-factor if and only if the source 3DM instance
has a perfect matching.
\end{theorem}

\begin{proof}
Given a perfect source matching, put its gadgets in ON mode and all others
in OFF mode. The local lemma gives the required degree at every private
vertex and every center. Each source vertex is in exactly one selected
triple, and hence is incident with exactly one selected port.

Conversely, let $F$ be an $f$-factor. The local lemma assigns each gadget a
unique mode. Let $M'$ consist of the source triples whose gadgets are ON.
An original source vertex is incident only with port edges and has required
degree one. Since an ON gadget selects all three of its ports and an OFF
gadget selects none, exactly one triple of $M'$ contains each source vertex.
Thus $M'$ is a perfect 3-dimensional matching.
\end{proof}

The construction has seven private vertices and ten edges per source triple,
so it is polynomial. Membership in NP follows by checking a proposed edge
subset and all incidence degrees. Together with balancedness, this proves
NP-hardness and NP-completeness for rank at most four.

\begin{corollary}[Perfect $f$-matching]
The same instances give NP-completeness for perfect $f$-matching.
\end{corollary}

\begin{proof}
Every target edge contains a vertex with $f$-value one: a helper on each
$U$-edge, a helper or $r$-vertex on each $B$-edge, an $r$-vertex on $A_e$,
and a source vertex on each port. Consequently, in any perfect $f$-matching
$x$, every edge multiplicity satisfies $x(e)\le1$. Thus $x$ is a 0/1
vector and is exactly an $f$-factor.
\end{proof}

\section{Bounded degree and the rank dichotomy}

For the bounded-occurrence source, the target degrees are
\[
\deg(c_e)=4+2+1+3=10,
\quad \deg(h_{e,i})=2,
\quad \deg(r_{e,i})=2,
\]
and every source vertex has degree at most three. Therefore
$\Delta(H)\le10$ and $\max_v f(v)=5$. This proves the strengthened form of
Theorem 1.

Combining the present theorem with the prior rank-at-most-three algorithm
and rank-five hardness gives the precise rank-bound classification:
\begin{center}
\begin{tabular}{ccl}
\toprule
rank bound & $f$-factor and perfect $f$-matching & status\\
\midrule
$r\le3$ & on balanced hypergraphs of rank at most $r$ & polynomial\\
$r\ge4$ & already on balanced hypergraphs of rank at most $r$ & NP-complete\\
\bottomrule
\end{tabular}
\end{center}
The second line means that the class with rank bound $r$ contains a hard
rank-four subclass; it does not assert that every input has rank at least
four. The construction is non-uniform.

\section{Mengerian perfect matching}

We use the vertex expansion operation from the prior literature, not as a
new notion. The next proposition records the required details.

\begin{proposition}[Expansion equivalence]
For any hypergraph $H$ and $f:V(H)\to\Z_{\ge0}$, $H$ has a perfect
$f$-matching if and only if $H^f$ has an ordinary perfect matching. Moreover,
$\operatorname{rank}(H^f)\le\operatorname{rank}(H)$.
\end{proposition}

\begin{proof}
If $x$ is a perfect $f$-matching, then for every $v$ the incident
occurrences of original edges number exactly $f(v)$. Assign the $f(v)$
clones of $v$ bijectively to these occurrences. Choosing the assigned clone
at every endpoint of each occurrence produces expanded edges that are
pairwise disjoint and cover all clones. Conversely, project a perfect
matching of $H^f$ to original edges and let $x(e)$ be the number of matching
edges projecting to $e$. Counting the clones of each $v$ gives
$\sum_{e\ni v}x(e)=f(v)$. Expansion replaces, rather than adds, one vertex
at each endpoint, so rank cannot increase.
\end{proof}

\begin{proposition}
If $H$ is balanced, then $H^f$ is Mengerian.
\end{proposition}

\begin{proof}
Balanced hypergraphs are Mengerian. Let $g$ be any expansion vector on the
clones of $H^f$. If the clones of $v$ in $H^f$ are
$v_1,\ldots,v_{f(v)}$, define
\[
h(v)=\sum_{i=1}^{f(v)}g(v_i).
\]
Expanding in two stages gives a natural isomorphism
$(H^f)^g\cong H^h$: the final clone of $v$ is identified by the pair
$(i,j)$ consisting of an intermediate clone and a clone inside it. Hence
matching and cover numbers agree, and
\[
\nu((H^f)^g)=\nu(H^h)=\tau(H^h)=\tau((H^f)^g).
\]
Thus $H^f$ is Mengerian.
\end{proof}

For the bounded-occurrence construction, every edge contains exactly one
center, and every center has $f$-value five. Hence
\[
|E(H^f)|=5|E(H)|=50|M|,
\qquad |V(H^f)|=3n+11|M|.
\]
The expanded center clones have degree ten. Each helper and each $r$-vertex
has degree $5\cdot2=10$, while each source vertex has degree at most
$5\cdot3=15$. Consequently $\Delta(H^f)\le15$, and rank remains at most
four.

\begin{proof}[Proof of Theorem 2]
Reduce from bounded-occurrence 3DM-3, construct $H$ as above, and output
$H^f$. The preceding propositions show that the output is Mengerian, has
rank at most four and maximum degree at most fifteen, and has a perfect
matching exactly when the source instance is positive. Perfect Matching is
in NP because a selected edge subset can be checked for pairwise
disjointness and full vertex coverage in polynomial time.
\end{proof}

\begin{corollary}[Ideal hypergraphs]
Perfect Matching is NP-complete for ideal hypergraphs of rank at most four,
even with maximum degree at most fifteen.
\end{corollary}

\begin{proof}
Every Mengerian hypergraph is ideal, so this is an immediate consequence of
Theorem 2.
\end{proof}

\section{A protected selector interpretation}

The gadget can be viewed as a rank-reduction replacement for the selector in
the earlier rank-five construction. Under the protected attachment pattern
used here, the old selector has two local behaviors: four fallbacks with no
external incidence, or one large selector with three external incidences.
The subsystem $B_1,B_2,A$ has the same external behavior. The equations
$b_1+a=b_2+a=1$ synchronize the two rank-four selector edges, and the center
requirement five forces the port sum to be either zero or three. The
balancedness proof is specific to the stated private incidence pattern; no
universal rank-reduction theorem is claimed.

\section{Independent verification}

The proof is independent of computation. As a separate implementation
check, exact integer programs enumerated all 256 subsets of the eight
possible triples on a $2\times2\times2$ source universe. The source and
target answers agreed, direct strong-cycle enumeration found no odd strong
cycle, and an independent odd-square incidence-matrix test found no
forbidden submatrix in the tested one- and two-triple cases. The expansion
solver agreed with the original f-factor solver on all 256 instances; all
193 instances satisfying the occurrence-at-most-three condition obeyed the
degree bounds, and all 120 perfect covers of the one-triple expansion
projected to valid multiplicity vectors. These checks are supplementary
falsification evidence, not a proof of any theorem.

\section{Conclusion}

The rank-four gap for non-uniform balanced hypergraph $f$-factors and
perfect $f$-matchings is closed by a rank-four balanced selector. The same
construction, followed by rank-preserving vertex expansion, gives bounded
rank hardness for Mengerian and ideal hypergraphs. We deliberately leave
uniform rank-four and counting variants outside the scope of this paper.

\bibliographystyle{plain}
\bibliography{references}
\end{document}
